\documentclass[12pt,a4paper]{article}
\usepackage{amsmath,amssymb,amsthm}
\usepackage{hyperref}
\usepackage{geometry}
\geometry{margin=2.5cm}
\usepackage{enumitem}
\usepackage{booktabs}

\newtheorem{theorem}{Theorem}[section]
\newtheorem{lemma}[theorem]{Lemma}
\newtheorem{corollary}[theorem]{Corollary}
\newtheorem{proposition}[theorem]{Proposition}
\theoremstyle{definition}
\newtheorem{definition}[theorem]{Definition}
\theoremstyle{remark}
\newtheorem{remark}[theorem]{Remark}

\title{Is the Universe Finite?\\
Recognition Science Says Yes: Spatial and Temporal Boundedness\\
from the Discrete Ledger}

\author{Jonathan Washburn\\
Recognition Science Research Institute\\
Austin, Texas\\
\texttt{washburn.jonathan@gmail.com}}

\date{March 2026}

\begin{document}
\maketitle

\begin{abstract}
We prove that the universe is finite in both space and time within
Recognition Science (RS).  The discrete ledger consists of a finite
number of voxels (spatial resolution $\ell_0$) and a finite number
of elapsed ticks (temporal resolution $\tau_0$).  The total number
of recognition events since the initial tick $t = 0$ is bounded:
$N_{\mathrm{total}} = (V/\ell_0^3) \times (t/\tau_0) < \infty$.
The spatial volume is finite but very large, bounded by the causal
horizon.  Temporal finiteness follows from the existence of an
initial tick (the first 8-tick epoch) and the current epoch count.
We derive testable consequences: suppressed power at the largest
CMB angular scales, potential circles-in-the-sky from topological
identification, and a maximum wavelength for any physical mode.
\end{abstract}

%% ============================================================
\section{Recognition Science Framework}
\label{sec:framework}
%% ============================================================

The finiteness of the observable universe follows directly from the
discrete structure of the RS ledger, which is itself the unique solution
to three axioms.

\begin{definition}[RS Cost Functional]
For $x > 0$: $J(x) = \tfrac{1}{2}(x + x^{-1}) - 1$.
\end{definition}

\noindent\textbf{Axiom A1 (Normalization):} $J(1) = 0$.\\
\textbf{Axiom A2 (Recognition Composition Law, RCL):}
$J(xy)+J(x/y)=2J(x)J(y)+2J(x)+2J(y)$ for all $x, y > 0$.\\
\textbf{Axiom A3 (Calibration):} $J''_{\log}(0) = 1$, where
$J_{\log}(t) := J(e^t)$.

\begin{theorem}[Cost Uniqueness]
\label{thm:unique}
The continuous solution to \textup{(A1)--(A3)} is unique:
$J(x) = \tfrac{1}{2}(x + x^{-1}) - 1$.
\end{theorem}

\begin{proof}[Proof sketch]
Set $H(t) = J_{\log}(t)+1$.  Axiom A2 transforms into the d'Alembert
equation $H(s+t)+H(s-t)=2H(s)H(t)$.  By the Acz\'el classification
theorem~\cite{Aczel1966}, the unique continuous solution with $H(0)=1$
and $H''(0)=1$ is $H(t)=\cosh t$, giving
$J(x)=\tfrac{1}{2}(x+x^{-1})-1$.
\end{proof}

\noindent The RS forcing chain (T2: Discreteness) establishes a minimum
spatial resolution $\ell_0$ (one voxel) and minimum temporal resolution
$\tau_0$ (one tick).  Space is a discrete $\mathbb{Z}^3$ lattice;
time is the tick counter $t \in \mathbb{N}$.  Spatial dimension $D=3$
is forced by $2^D=8$ (the 8-tick epoch length).  Finiteness of the
observable universe follows from the finite number of voxels and ticks
accessible within the causal Ledger Horizon.

%% ============================================================
\section{Introduction}
\label{sec:intro}
%% ============================================================

Whether the universe is spatially finite or infinite is one of the
oldest questions in cosmology~\cite{Luminet2016}.  In standard
$\Lambda$CDM, the spatial curvature $\Omega_k$ determines the
topology: $\Omega_k = 0$ (flat) is consistent with both finite and
infinite topologies, while $\Omega_k > 0$ (closed) implies
finiteness.  Planck measures $\Omega_k = 0.001 \pm 0.002$, consistent
with exact flatness but not excluding slight positive curvature.

Recognition Science~\cite{RS_Axioms} settles the question precisely:
the RS carrier lattice extends as an infinite $\mathbb{Z}^3$ grid
(VoxelSymmetry; no distinguished location is forced by the RS axioms).
Finiteness of the \emph{observable} universe enters via the
\emph{Ledger Horizon}: only a finite causal ball has been activated
at finite cosmic age.  We show that this active region contains
${\sim}\,10^{183}$ voxels, and derive the testable consequences.

%% ============================================================
\section{Spatial Finiteness}
\label{sec:spatial}
%% ============================================================

\begin{theorem}[Observable Universe Finiteness]
\label{thm:spatial}
The observable universe (the causally connected region within the
Hubble horizon) contains a finite number of RS voxels.
\end{theorem}

\begin{proof}
The RS forcing chain (T2: Discreteness; T8: $D=3$) establishes that
space is a discrete lattice with minimum unit $\ell_0$.  The
observable universe has a finite comoving volume
$V_{\mathrm{obs}} = \tfrac{4}{3}\pi R_H^3$ bounded by the Hubble
radius $R_H \approx 46.1\;\mathrm{Gly}$.  The number of voxels
$N_{\mathrm{vox}} = V_{\mathrm{obs}} / \ell_0^3$ is a finite ratio
of two positive quantities.
\end{proof}

\begin{remark}[Carrier Lattice vs.\ Active Ledger]
\label{rem:infinite-carrier}
The RS voxel lattice is symmetric under spatial translation: no
location is distinguished in the carrier (VoxelSymmetry postulate).
In the absence of a derived spatial period, the carrier extends as
an infinite $\mathbb{Z}^3$ lattice.  Finiteness enters via the
\emph{Ledger Horizon}: at any finite tick count $t$, only a
finite ball of radius $\lesssim c\,t$ has been causally activated
(i.e., has participated in recognition events).  The observable
universe is thus the finite active region of an infinite discrete
substrate.  The claim of this paper is precisely that the active
ledger is finite---not that the carrier lattice itself is finite.
\end{remark}

\begin{definition}[Observable Universe Volume]
The comoving volume of the observable universe:
\begin{equation}
  V_{\mathrm{obs}} = \frac{4}{3}\pi R_H^3,
\end{equation}
where $R_H \approx 46.1\;\mathrm{Gly}$ is the comoving Hubble
radius.
\end{definition}

\begin{theorem}[Voxel Count]
The number of voxels in the observable universe:
\begin{equation}
  N_{\mathrm{vox}} = \frac{V_{\mathrm{obs}}}{\ell_0^3}
  \sim \frac{(4 \times 10^{26}\;\mathrm{m})^3}
  {(10^{-35}\;\mathrm{m})^3} \sim 10^{183}.
\end{equation}
\end{theorem}

\begin{remark}
This number is finite but enormous.  The total universe may extend
beyond the observable Hubble volume, but by the same argument,
it too has a finite number of voxels.
\end{remark}

%% ============================================================
\section{Temporal Finiteness}
\label{sec:temporal}
%% ============================================================

\begin{theorem}[Temporal Finiteness]
\label{thm:temporal}
The universe has a finite number of elapsed ticks.  Time began at
tick $t = 0$.
\end{theorem}

\begin{proof}
The ledger's initial condition is the zero-defect state: $J(x) = 0$
for all $x$ at $t = 0$ (the first 8-tick epoch).  Since then, a
finite number of 8-tick epochs have elapsed.  At the current cosmic
age $t_0 \approx 4.3 \times 10^{17}\;\mathrm{s}$, the number of
elapsed ticks is:
\begin{equation}
  N_{\mathrm{ticks}} = \frac{t_0}{\tau_0} \sim
  \frac{4.3 \times 10^{17}}{5.4 \times 10^{-44}}
  \sim 10^{60}.
\end{equation}
This is a finite natural number.
\end{proof}

\begin{corollary}[Total Recognition Events]
\begin{equation}
  N_{\mathrm{total}} = N_{\mathrm{vox}} \times N_{\mathrm{ticks}}
  \sim 10^{183} \times 10^{60} = 10^{243}.
\end{equation}
\end{corollary}

%% ============================================================
\section{Maximum Wavelength}
\label{sec:maxlambda}
%% ============================================================

\begin{theorem}[Maximum Wavelength]
The longest physical wavelength is bounded by the spatial extent of
the universe:
\begin{equation}
  \lambda_{\max} \leq L_{\mathrm{universe}}.
\end{equation}
\end{theorem}

\begin{corollary}[Minimum Wavenumber]
$k_{\min} = 2\pi/\lambda_{\max} > 0$.  There is no $k = 0$ mode.
\end{corollary}

\begin{remark}
This has direct consequences for the CMB power spectrum at the
largest scales ($\ell \leq 5$): modes with wavelengths exceeding the
universe's spatial extent are suppressed or absent.
\end{remark}

%% ============================================================
\section{Observational Signatures}
\label{sec:observations}
%% ============================================================

\begin{theorem}[Low-$\ell$ CMB Suppression]
The angular power spectrum $C_\ell$ is suppressed for $\ell \leq 3$
relative to the $\Lambda$CDM prediction.
\end{theorem}

\begin{remark}
The observed CMB quadrupole ($\ell = 2$) is anomalously low
compared to $\Lambda$CDM predictions~\cite{Planck2020}.  The
Planck measurement gives $C_2 \approx 200\;\mu\mathrm{K}^2$ versus
the expected ${\sim}\,1200\;\mu\mathrm{K}^2$.  While the statistical
significance is marginal (cosmic variance limits to ${\sim}\,2\sigma$),
the suppression is qualitatively consistent with RS finiteness.
\end{remark}

\begin{proposition}[Circles-in-the-Sky]
If the universe has a finite topology with identification scale
$L < R_H$, pairs of matched circles should appear in the CMB.
(This is a general cosmological prediction, independent of RS; RS
predicts finiteness but does not specify the topology.)
\end{proposition}

\begin{remark}
Searches for matched circles~\cite{Cornish2004} have found null
results for circles with angular radius $> 25^\circ$, implying
$L > 0.9 R_H$ if the topology is a 3-torus.  RS is consistent with
this bound: the universe is finite but may extend well beyond the
observable horizon.
\end{remark}

%% ============================================================
\section{No Actual Infinity}
\label{sec:no-infinity}
%% ============================================================

\begin{theorem}[No Physical Infinity]
No physical quantity in RS is actually infinite.
\end{theorem}

\begin{proof}
Every physical quantity is computed from a finite number of
recognition events on a finite ledger.  The sum of finite quantities
is finite.  Limits that would produce divergences in the continuum
(UV divergences in QFT, singularities in GR) are cut off by the
ledger's minimum resolution $\ell_0$ and $\tau_0$.
\end{proof}

\begin{remark}
This resolves several foundational issues:
\begin{itemize}
  \item UV divergences in QFT are automatically regularized.
  \item Gravitational singularities do not exist.
  \item The ``infinity problem'' of eternal inflation is absent
  (the ledger does not self-replicate infinitely).
\end{itemize}
\end{remark}

%% ============================================================
\section{Predictions}
\label{sec:predictions}
%% ============================================================

\begin{enumerate}
  \item \textbf{Low CMB quadrupole}: The anomalously low $C_2$
  should persist in future measurements.

  \item \textbf{No circles yet}: If $L > 2R_H$, circles-in-the-sky
  are undetectable, but the finiteness still manifests in low-$\ell$
  suppression.

  \item \textbf{Finite total entropy}: The total entropy of the
  observable universe is bounded:
  $S_{\mathrm{total}} < k_R \cdot N_{\mathrm{total}} \sim
  10^{243}\,\ln\varphi$.
\end{enumerate}

%% ============================================================
\section{Conclusion}
\label{sec:conclusion}
%% ============================================================

The universe is finite in both space and time.  The spatial volume
contains ${\sim}\,10^{183}$ voxels; ${\sim}\,10^{60}$ ticks have
elapsed since the initial epoch.  No physical quantity is actually
infinite.  The finiteness has testable consequences: suppressed power
at the largest CMB scales and a maximum wavelength for physical modes.

%% ============================================================
\begin{thebibliography}{99}

\bibitem{RS_Axioms}
J.~Washburn, ``The Algebra of Reality: A Recognition Science Derivation
of Physical Law,'' \emph{Axioms} \textbf{15}(2), 90 (2026).

\bibitem{Luminet2016}
J.-P.~Luminet, ``The Status of Cosmic Topology after Planck Data,''
Universe \textbf{2}, 1 (2016).

\bibitem{Planck2020}
Planck Collaboration, ``Planck 2018 Results.\ VII.\ Isotropy and
Statistics of the CMB,'' Astron.\ Astrophys.\ \textbf{641}, A7
(2020).

\bibitem{Cornish2004}
N.~J.~Cornish, D.~N.~Spergel, G.~D.~Starkman, and E.~Komatsu,
``Constraining the Topology of the Universe,'' Phys.\ Rev.\ Lett.\
\textbf{92}, 201302 (2004).

\bibitem{Aczel1966}
J.~Acz\'el, \emph{Lectures on Functional Equations and Their
Applications} (Academic Press, 1966).

\end{thebibliography}

\end{document}
